% Beamer's default font size is 11 points. It is possible to set the
% default font size to any of 8, 9, 10, 11, 12, 14, 17, 20
\documentclass[13pt]{beamer}

\usetheme{Warsaw}
%\usecolortheme{dolphin}
\usecolortheme{seahorse}
%\useinnertheme{circles}
\usefonttheme{professionalfonts}
%\usefonttheme{structurebold}
%\usefonttheme{default}

\setbeamertemplate{navigation symbols}{}
 
\usepackage[english]{babel}
\usepackage[latin1]{inputenc}
\usepackage{times}
\usepackage[T1]{fontenc}

% These things can change:
\newcommand{\baseurl}{http://cell.vtt.fi/latex}
\urlstyle{sf}
\newcommand{\plus}{\makebox[2em]{+}}
\newcommand{\minus}{\makebox[2em]{-}}


\begin{document}

\logo{\includegraphics[width=1.4cm]{img/vttplain}\quad}
\title{\scalebox{1.5}{\textrm{\LaTeX}} Course 2011}
\subtitle{$\mathcal{F}$rom \scalebox{1.2}{$\beta$}asics to
\scalebox{1.2}{$\alpha$}dvanced
\scalebox{1.2}{$\tau$}ypesetting}
\author{Arho Virkki}
\institute{\textsc{VTT Technical Research Centre of Finland}}
\date{\scalebox{0.8}{\textcolor{gray}{2011-10-04}}}

\begin{frame}
  \titlepage
\end{frame}


\begin{frame}\frametitle{About the Course}
\begin{block}{Course Page:}
\url{\baseurl}
\end{block}\bigskip

Suggested preliminary knowledge:\smallskip

\framebox[5em][l]{\textbf{original:}} 
\url{http://tobi.oetiker.ch/lshort/lshort.pdf}\\
\framebox[5em][l]{\textbf{copy:}} \url{\baseurl/extras/lshort.pdf}\bigskip

Suomeksi:\smallskip

\hspace*{-1em}\framebox[4em][l]{\textbf{alkup.:}}
\url{http://www.tex.ac.uk/CTAN/info/lshort/finnish/lyhyt2e.pdf}\\
\hspace*{-1em}\framebox[4em][l]{\textbf{kopio:}}
\url{\baseurl/extras/lyhyt2e.pdf} \bigskip

\end{frame}


\begin{frame}\frametitle{About the Course\dots}
Suggested book: \\
\textbf{Helmut Kopka and Patrick W. Daly, Guide to LaTeX,
4th edition.} \\
In my opinion, the book is very thoroughly written and works 
perfectly as a reference.\bigskip

Other suggested books: \\
\textbf{The LaTeX Companion, The LaTeX Web Companion, The LaTeX Graphics
Companion} \\
These are the fundamental and also the most often cited books. Their only
disadvantage is that there are even more pages to read through. 
\end{frame}


\begin{frame}\frametitle{About the Course\dots}
There appears not to be both a \textbf{complete} and a \textbf{concise}
electronic reference to the system. In contrast, the amount of documentation
is huge, and most of it is more or less
outdated.\bigskip

$\to$ The book is a valuable investment.\bigskip

Some of the electronic references are rather complete: teTeX 
hypertext help is rather complete, and the TeXlipse plug-in for Eclipse knows
quite the well the basic commands.
\end{frame}

\begin{frame}\frametitle{About the Course\dots}
Most of the texts in natural sciences are written with \LaTeX{} 
nowadays.\bigskip

\LaTeX{} is actually part of the scientific heritage, including the
other open sourced tools like Linux, PostgreSQL, the BSD-system, Apache,
etc. 
\bigskip

\emph{Keywords:} markup languages, GNU and open source.
\end{frame}


\begin{frame}\frametitle{What is \LaTeX{}?}
\LaTeX{}, or precisely \LaTeXe{}, is a typesetting system for structured
documents, including

\begin{itemize}
  \item books
  \item scientific papers
  \item handouts
  \item research reports
  \item slides --- like now.
\end{itemize}
\end{frame}


\begin{frame}\frametitle{What is \LaTeX{}?\dots}
The document format is structured, which means that there are
\begin{itemize}
  \item logical and (preferably only few) %(\verb|\begin{document}|) ja
  \item typographical %(\verb|\textbf{lihavoituna}|)
\end{itemize}
commands mixed with the manuscript text\footnote{The manuscript  has the role
of ``source code'' in this context.}.\bigskip

The text-formatted source can be compiled into
different formats, the most common being
\begin{itemize}
  \item PDF (Portable Document Format) and
  \item HTML (HyperText Markup Language).
\end{itemize}
\end{frame}


\begin{frame}\frametitle{This course is for you, if you\ldots}

This course is for you, if you\ldots
\begin{itemize}
  \item compose books,
  \item publish in natural sciences,
  \item compose mathematical texts,
  \item generate (pdf) documents automatically from software,
  \item just like software that (just) works, or
  \item just like free / open source software.
\end{itemize}
\end{frame}


\begin{frame}\frametitle{Course contents}

Course contents
\begin{enumerate}
  \item short introduction
  \item logical constructs: lists, environments,\dots
  \item typographical building blocs: lines, boxes, fonts,\ldots
  \item mathematical typesetting: 
  $\displaystyle\oint_C \frac{3\pi}{\sqrt[n]{\frac{z-1}{z+1}}}\, \mathrm{dz}$,
  \dots
  \item graphics: adding and drawing images,\dots
  \item other tools: BibTeX for references, drawing images,\dots
\end{enumerate}
\end{frame}


\begin{frame}\frametitle{Course contents}
\begin{center}
\scalebox{3}{\textcolor{gray}{Welcome}}
\rotatebox{-45}{\scalebox{2}{\textsl{to}}}
\rotatebox{45}{\scalebox{2}{\textrm{the}}}
\raisebox{1em}{\rotatebox{-180}{\scalebox{2}{\textsc{Course!}}}}%
\scalebox{2}{\textsc{!}}

\begin{displaymath}
X\beta = y + \epsilon \quad \Leftrightarrow \quad
\begin{bmatrix}
1 & x_{11} &  \cdots & x_{1M} \\
1 & x_{21} &  \cdots & x_{1M} \\
\vdots &  &  & \vdots \\
\\
\vdots &  &  & \vdots \\
1 & x_{N1} & \cdots & x_{NM} \\
\end{bmatrix}
\raisebox{1.5em}{$\begin{bmatrix}
\beta_0 \\ \beta_1 \\ \vdots \\ \beta_M 
\end{bmatrix}$} = 
\begin{bmatrix}
y_1 \\ y_2 \\ \vdots \\ \\ \vdots \\ y_N 
\end{bmatrix} + 
\begin{bmatrix}
\epsilon_1 \\ \epsilon_2 \\ \vdots \\ \\ \vdots \\ \epsilon_N 
\end{bmatrix}
\end{displaymath}
\end{center}
\end{frame}


\end{document}
